% Generated draft: requires mathematical and source review.
\documentclass[11pt]{article}
\usepackage[margin=25mm]{geometry}
\usepackage{fontspec}
\setmainfont{Latin Modern Roman}
\usepackage{amsmath,amssymb,mathtools,unicode-math}
\setmathfont{Latin Modern Math}
\usepackage{enumitem,needspace,xurl,hyperref}
\setlength{\parindent}{0pt}
\setlength{\parskip}{5pt}
\setlength{\emergencystretch}{3em}
\newcommand{\N}{\mathbb{N}}
\newcommand{\Z}{\mathbb{Z}}
\newcommand{\Q}{\mathbb{Q}}
\newcommand{\R}{\mathbb{R}}
\newcommand{\C}{\mathbb{C}}
\newcommand{\F}{\mathbb{F}}
\newcommand{\A}{\mathbb{A}}
\newcommand{\PP}{\mathbb{P}}
\newcommand{\E}{\mathbb{E}}
\newcommand{\eps}{\varepsilon}
\newcommand{\dd}{\,\mathrm{d}}
\begin{document}
\section*{UnsolvedMath ID: 30003949}
\textbf{Source problem code:} catalogue entry
\subsection{Definitions and mathematical statement}
A hyperbolic surface is a smooth surface \(S\) equipped with a Riemannian metric \(g\) of constant Gaussian curvature \(-1\). For functions on \(S\), use the nonnegative Laplace--Beltrami convention
\[
\Delta_g f=-\operatorname{div}_g(\nabla_g f).
\]
If \(S\) is compact and without boundary, its Laplace spectrum can be listed, with multiplicities, as
\[
0=\lambda_0(S,g)\leq\lambda_1(S,g)\leq\lambda_2(S,g)\leq\cdots,
\qquad \lambda_j(S,g)\longrightarrow+\infty.
\]
This spectral notation is conditional on the absence of boundary; the supplied problem does not specify whether boundary is allowed. If boundary is allowed, a self-adjoint boundary condition must be specified before the eigenvalues are determined.

\textbf{Problem (supplied wording).} Construct a McShane-type identity involving the eigenvalues of the Laplacian of a compact hyperbolic surface.

The supplied record does not define what qualifies as a McShane-type identity in this spectral setting. In particular, it gives no summation index set, summand, convergence requirement, or prescribed value for an identity. It also does not specify whether the identity is to hold for every compact hyperbolic surface, for a restricted class, or under additional hypotheses. Thus the requested construction is an open-ended research problem; a fully quantified identity cannot be extracted from the supplied statement without additional choices.
\subsection{Short English statement}
The problem asks for a McShane-type identity expressed using Laplace eigenvalues of a compact hyperbolic surface. The intended form of the identity, its range of validity, and any boundary conditions remain unspecified. The supplied open-status label is dated curation and has not been independently verified.
\subsection{Sources}
\begin{itemize}
\item Local UnsolvedMath registry, permanent ID \(30003949\), title ``McShane Identities for Laplace Spectra,'' source problem code \texttt{catalogue entry}, in the Oberwolfach Reports Open Problems collection. The record marks the statement as exact and publication-ready; these are supplied editorial qualifications, and the formulation issues described above remain unresolved.
\item Supplied source report: \emph{New Trends in Teichmüller Theory and Mapping Class Groups} (2019), \emph{Oberwolfach Reports}, pp.~2475--2534, DOI \texttt{10.4171/owr/2018/40}. \url{https://doi.org/10.4171/owr/2018/40}. The record gives proposed year 2018. Its literature-source label identifies this as ``Oberwolfach Report 40, A'Campo problem''; the proposer field itself is unspecified.
\item Supplied literature source: ``McShane-type identity for closed surfaces.'' \url{https://arxiv.org/abs/1203.3860}.
\item Supplied literature assessment, checked on exactly 2026-08-21: status ``open.'' The assessment reports known McShane-type identities for geodesic lengths and higher-rank representations, and reports finding no identity in terms of Laplacian eigenvalues answering this problem. This is supplied dated curation derived from the OWR material, not an independently verified literature assessment. The cited sources and time-sensitive status require independent verification.
\end{itemize}
\end{document}
